\begin{lemma}[Missing colors are excluded from an elementary tail]
Let $(T,\varphi_n)$ be a positive-rung extended Tashkinov tree in a
finite loopless multigraph with individual edge identities, let
$\mathcal B$ be a hierarchy base under $\varphi_n$, and let $H$ be a
hierarchy on this base.  Put
\[
 q=H.\mathit{dividerCount},\qquad
 A=\noderef{HierarchyStarredStage}(H,q),\qquad
 \ell=|T|-|A|,
\]
and, for a natural number $j$, let
\[
 T(y_j)=\noderef{TreeSequencePrefix}
       \bigl(T,\min\{|A|+j,|T|\}\bigr).
\]
Let $n$ be the terminal rung index and let $T_n$ be the terminal ladder
tree.  Assume that $V(T(y_{\ell-1}))$ is elementary under $\varphi_n$.

Fix $i$ with $1\leq i<\ell$, and suppose that a color $\beta$ is missing
at the exact last vertex $y_i$ of $T(y_i)$.  If $t$ is an edge position
satisfying
\[
             |A|\leq t<|A|+i
 \quad\hbox{and}\quad
             \varphi_n(T(y_i).\mathit{edge}(t))=\beta,
\]
then all of the following exceptional conclusions hold:
\begin{enumerate}
\item $q=0$, $A=\mathcal B.\mathit{core}=T_n$, and $t=|A|$;
\item the actual certified step at index $n$ is an SE or RE constructor;
\item for an individual edge identity $f_n$, one has
\[
 F_n=\{f_n\},\qquad
 T(y_i).\mathit{edge}(t)=f_n,\qquad
 \varphi_n(f_n)=\beta\in D_n,
\]
where $D_n=\noderef{TashkinovStoredColors}(\mathcal T,n)$.
\end{enumerate}
Thus no $\beta$-colored exact edge occurs from $A$ through $y_i$, apart
from the first displayed SE/RE connecting seed in precisely the stated
zero-divider case.  The assertion concerns edge identities, not endpoint
pairs, so it remains valid in the presence of parallel edges.
\end{lemma}
\begin{proof}
Write $B=T(y_{\ell-1})$.  The range $1\leq i<\ell$ gives
$i\leq\ell-1$.  The last-starred-stage definition and the nested,
core-inside-first, and terminal-stage fields of
\noderef{TashkinovHierarchy} show that $A$ is an exact ordered initial
segment of $T$.  Hence $|A|+\ell=|T|$, the minimum defining $T(y_i)$ is
$|A|+i$, and $T(y_i)$ is an exact ordered initial segment of $B$.
By \noderef{TreeSequencePrefix}, its last vertex is the entry of $T$ at
index $|A|+i$.  All vertex and edge equalities used below are the
entrywise equalities of \noderef{TreeSequenceInitialSegment}; no equality
of endpoint pairs is substituted for an edge identity.

Let $n=\operatorname{Fin.last}(\mathit{ett.rounds})$.  Since the ETT has
positive rung, $n$ is a positive series index.  Inspect the actual
certificate $\mathcal T.\mathit{certified}(n)$ supplied by
\noderef{TashkinovSeries}.  The initial constructor of
\noderef{TashkinovSeriesStep} is impossible at this positive index, so
the certificate is PE, SE, or RE.  The generating-coloring field of
\noderef{ExtendedTashkinovTree} identifies its current coloring with
$\varphi_n$.

We next locate the certified Tashkinov-augmentation route.  In the PE
case, the constructor supplies a
$\varphi_n$-\noderef{TashkinovClosure} from $T_n$ to the final certified
output.  Thus it supplies a
$\varphi_n$-\noderef{TashkinovAugmentationSequence} from $T_n$ to that
output.  The PE clause of \noderef{HierarchyBase} makes $T_n$ an exact
ordered initial segment of the base core, and the hierarchy places the
core inside $A$, then $T(y_i)$, then $T$, and finally the certified
output.  Applying
\noderef{TashkinovAugmentationSequenceRestrict} at these exact intermediate
prefixes gives the route through $A$ and $T(y_i)$.

In an SE or RE constructor, the certificate instead supplies an
individual edge $f_n$, the singleton equation $F_n=\{f_n\}$, the exact
one-edge seed $S=T_n+f_n$ in the sense of
\noderef{TreeSequenceAdjoinEdge}, and a $\varphi_n$-Tashkinov closure from
$S$ to the final output.  If $q>0$, the exact conclusion of
\noderef{HierarchyNonparallelFirstEdge} places this same identity $f_n$
at position $|T_n|$ of the first positive hierarchy stage.  The seed and
that stage are prefixes of the same certified output and the seed has
length $|T_n|+1$, so their entrywise equations make $S$ a prefix of the
first stage and hence of $A$.  Restriction again gives the augmentation
route through $A$ and $T(y_i)$.

If $q=0$ in the SE/RE case, the nonparallel clause of
\noderef{HierarchyBase} gives
\[
                              A=T_n.                           \tag{1}
\]
Here the only displayed edge immediately after $A$ that is not itself a
Tashkinov-augmentation edge is the adjoined seed $f_n$.  Every later
edge lies in the closure route from $S$ and is an augmentation edge.
For completeness, this last assertion follows directly by induction on
the finite augmentation chain.  Every step increases the length by one
and retains every earlier entry.  If the chain starts at a length at
most $t$ and ends at $T(y_i)$, whose length is greater than $t$, take
the first member whose length is greater than $t$.  Its predecessor has
length at most $t$, and the one-step length equation forces the two
lengths to be $t$ and $t+1$.  The new last edge of that step is therefore
the exact displayed identity at position $t$.
We have therefore proved the following exact dichotomy for the position
$t$ in the statement:
\[
 \begin{array}{l}
 \text{either the edge at }t\text{ is the new last edge of a certified }
 \varphi_n\text{-Tashkinov augmentation,}\quad\text{or}\\[1mm]
 q=0,\ A=T_n,\ t=|A|,\ \Theta_n\in\{\mathrm{SE},\mathrm{RE}\},
 \ F_n=\{f_n\},\text{ and the edge at }t\text{ is }f_n.
 \end{array}                                                   \tag{2}
\]
The second line uses \noderef{HierarchyNonparallelFirstEdge} with the
zero-divider terminal-stage identity, so its equality is the exact
$e_1=f_n$ equality required in the manuscript.

We rule out the first line of (2).  Expand the corresponding one-step
\noderef{TashkinovAugmentation}.  Its predecessor $X$ has length $t$,
its new last edge is the exact edge at position $t$, and its
$\varphi_n$-color is missing at some vertex $z$ of $X$.  Under the
hypothesis that this edge has color $\beta$, we have
\[
             \beta\in\overline\varphi_n(z).                   \tag{3}
\]
The predecessor is an exact prefix of $T(y_i)$, so $z$ occurs at an
index at most $t$.  The assumed strict inequality $t<|A|+i$ and vertex
injectivity in \noderef{TreeSequence} show that $z$ is not the last
vertex $y_i$, whose index is $|A|+i$.  Both $z$ and $y_i$ lie in
$V(T(y_i))\subseteq V(B)$.  Equation (3) and the assumed
$\beta\in\overline\varphi_n(y_i)$ now contradict the elementarity of
$B$ as defined by \noderef{ColoredVertexSetElementary}.  Therefore only
the exceptional second line of (2) can occur.  This proves $q=0$, (1),
$t=|A|$, the actual certified SE/RE alternative, $F_n=\{f_n\}$, the
exact edge identity, and $\varphi_n(f_n)=\beta$.

It remains to verify the asserted membership in $D_n$, rather than
merely label the seed exceptional.  In the SE constructor, the fields of
the actual \noderef{TashkinovSeriesStep} certificate and its
\noderef{MaximumDefectiveChoice} identify $f_n$ with the selected
defective edge, identify its color with the selected color
$\delta_n=\beta$, and put $\delta_n$ in the raw stored set $S_n$.  In the
RE constructor, the \noderef{RevisitingExtensionCertificate} return path
begins with the exact edge $f_n$ and alternates the certified colors
$\gamma_h,\delta_h$; hence
$\beta$ is one of these two colors, and the constructor puts both in
$S_n$.  Thus in either case
\[
                              \beta\in S_n.                    \tag{4}
\]

We claim that $\beta$ is not missing anywhere on $T_n$.  In the RE
constructor, the current coloring $\varphi_n$ is exactly the preceding
coloring and the constructor states that $T_n$ is closed under that
coloring.  The exact edge $f_n$ is a boundary edge of $T_n$ carrying
$\beta$.  The definition \noderef{ColoredVertexSetClosed} therefore
immediately excludes $\beta$ from the missing-color set of every vertex
of $T_n$.

For SE, write $\psi$ for the preceding coloring and $\pi=\varphi_n$ for
the selected coloring.  The constructor gives that $T_n$ is closed under
$\psi$, that $\pi$ is
$\bigl(T_n,D_{n-1},\psi\bigr)$-stable, that $\pi$ is proper, and that
$f_n$ is a boundary edge incident with the selected inside vertex $v$
and has $\pi(f_n)=\beta$.  Suppose that $\beta$ were missing somewhere
on $T_n$ under $\pi$.  The vertexwise clause of
\noderef{StableColoring} then makes $\beta$ missing somewhere on $T_n$
under $\psi$.  On the other hand, $\pi(f_n)=\beta$ and incidence with
$v$ show from \noderef{MissingColors} that $\beta$ is not missing at
$v$ under $\pi$; vertexwise stability says the same under $\psi$.
Thus some edge $g$ incident with $v$ has $\psi(g)=\beta$.  Because
$\beta$ is tree-missing under $\psi$, the incident-edge clause of
\noderef{StableColoring} fixes this exact edge, so $\pi(g)=\beta$.
Properness of $\pi$ from \noderef{IsProperEdgeColoring}, applied to the
two $\beta$-colored edges $f_n$ and $g$ incident with $v$, gives
$g=f_n$.  Hence $\psi(f_n)=\beta$.  This contradicts the
$\psi$-closedness of $T_n$, since $f_n$ is its boundary edge and
$\beta$ is missing somewhere on it.  This proves in SE as well that
\[
                    \beta\notin\overline\varphi_n(T_n).        \tag{5}
\]
Finally, \noderef{TashkinovStoredColors} defines $D_n$ as the union of
the raw stored sets through index $n$, with
$\overline\varphi_n(T_n)$ removed.  Combining (4) and (5) gives
$\beta\in D_n$.  This completes every clause of the exceptional
conclusion.
\end{proof}
